\chapter{Qué sigue}
\label{cap:siguientes}

\begin{enpocas}
Resumen de lo construido, cómo reportarlo en un artículo o reporte y experimentos que valen la
pena con FuegoNet.
\end{enpocas}

\section{Lo que ya tienes}

\begin{itemize}
  \item Una implementación propia, legible y sin dependencias externas de la
        arquitectura y la receta de entrenamiento de YOLO26: bloques C3k2, SPPF y
        C2PSA, cabeza sin DFL ni NMS, TAL con STAL, pérdida de dos cabezas con
        ProgLoss y optimizador MuSGD.
  \item Un modelo entrenado desde cero, medido con un protocolo propio y
        comparado contra la implementación de referencia.
  \item La red conectada a la misión de búsqueda (\cod{red:=fuegonet}) y exportable
        a ONNX.
\end{itemize}

\section{Cómo reportarlo}

\begin{enumerate}
  \item En el marco teórico, describe la arquitectura con la
        figura~\ref{fig:arquitectura} y la tabla~\ref{tab:capas}, y la asignación y la
        pérdida con las ecuaciones de los capítulos~\ref{cap:asignacion}
        y~\ref{cap:perdida}.
  \item En la metodología, da el dataset (tabla~\ref{tab:dataset}), los aumentos, los
        hiperparámetros del capítulo~\ref{cap:entrenar} y el protocolo de medición
        (mismo código para todas las redes).
  \item En los resultados, la tabla comparativa del capítulo~\ref{cap:evaluar} y las
        curvas de entrenamiento. Separa con claridad «arquitectura igual» de «pesos
        iniciales distintos».
  \item Cita YOLO26 por su documentación de Ultralytics, YOLOv10 por la idea de la
        doble asignación sin NMS, TOOD por TAL y el trabajo de Muon por el optimizador.
\end{enumerate}

\section{Experimentos sugeridos}

\begin{description}
  \item[Ablación] Quitar una pieza a la vez y medir: sin STAL (\cod{lado_min=0}),
        sin atención (\cod{atencion=False} en \cod{h22} y sin \cod{C2PSA}), SGD en
        lugar de MuSGD (\cod{--optimizador sgd}), sin ProgLoss ($w$ fijo en 0.5). Cada
        pieza debe justificar su lugar.
  \item[Escala] Entrenar FuegoNet-s (\cod{--escala s}, 10~M parámetros) y ver si el
        aumento de mAP compensa la latencia en la computadora a bordo.
  \item[Preentrenamiento propio] Entrenar primero con un dataset público de incendios
        reales (por ejemplo D-Fire) y después ajustar con el simulador. En el simulador
        FuegoNet ya supera al YOLO26n preentrenado, pero con imágenes reales el
        preentrenamiento suele pesar mucho más.
  \item[Varias semillas] Repetir FuegoNet-n y el YOLO26n desde cero con 3 semillas
        (\cod{--semilla}) y reportar media y desviación. Solo así se puede afirmar que
        una diferencia de 0.01 en mAP50-95 es real.
  \item[Fusionar BatchNorm] Plegar cada BatchNorm en su convolución y quitar la cabeza
        o2m antes de inferir, para acercar la latencia en CPU a la de Ultralytics.
  \item[Del simulador al mundo real] Probar la red con videos reales de incendios
        tomados con dron. Ahí se verá si los aumentos de color alcanzan o si hace
        falta aleatorizar más el simulador (texturas, luz, niebla).
  \item[Clase «distractor»] Agregar una tercera clase para lo que hoy da falsas
        alarmas (rocas oscuras, claros con sol) en lugar de dejarlo como fondo.
\end{description}

\begin{recuerda}
Una red que mide bien en validación todavía no está probada en vuelo. La prueba
final es la misión completa: cuánto tarda el Tarot 650 en encontrar y confirmar
el foco, y cuántas falsas alarmas da en el camino.
\end{recuerda}
